\chapter{Development Workflow \& Conventions}
\label{ch:workflow}

\section{Migrations: the one hard rule}

\textbf{Never hand-create a file inside \pathname{supabase/migrations}.} Always use the CLI:

\begin{lstlisting}[style=olkcodebash, language=bash, caption={Correct way to start a new migration}]
npx supabase migration new add_some_feature
\end{lstlisting}

This produces a correctly-timestamped file (\pathname{YYYYMMDDHHMMSS\_add\_some\_feature.sql}) in the exact chronological ordering the CLI expects. Migrations in this project are applied strictly in filename order — the timestamp \emph{is} the dependency graph. A hand-created file with a wrong or duplicate timestamp can silently apply out of order, or collide with a timestamp another migration already used, and either failure mode is much harder to debug than the two seconds it costs to run the command above.

\begin{olknote}
Look at the existing migration history in Chapter~\ref{ch:schema} for the house style: every non-trivial migration in this project opens with a plain-English comment block explaining \emph{why}, not just what — see \pathname{20260701182526\_rate\_limit\_messages\_and\_requests.sql} or \pathname{20260701182527\_wishlist\_fuzzy\_match\_and\_unique.sql} for examples worth imitating. A migration titled only with what it does (\emph{"add column bio"}) is worth a fraction of one that also records why it was needed.
\end{olknote}

\section{Keeping \pathname{schema.sql} honest}

\pathname{supabase/schema.sql} is a full point-in-time dump of the database, intended as a single-file, greppable reference. It is \textbf{not} automatically regenerated — it once drifted five migrations stale before being caught and regenerated (Chapter~\ref{ch:schema} tells that story). Nothing about the tooling prevents it happening again: regenerate it deliberately after a batch of migrations, with:

\begin{lstlisting}[style=olkcodebash, language=bash, caption={Regenerating the schema dump against the local database}]
npx supabase db dump --local --schema public > supabase/schema.sql
\end{lstlisting}

\begin{olkwarn}
Do not treat \pathname{schema.sql} as the source of truth if it disagrees with \pathname{supabase/migrations} — the migrations directory always wins, because it is what actually gets applied. Chapter~\ref{ch:schema}'s reference tables are built from the migrations, not from the dump, for exactly this reason. CI (below) does not check whether the dump is current — regenerating it after a batch of migrations is still a manual habit, not an enforced one.
\end{olkwarn}

\section{Linting and CI}

\begin{lstlisting}[style=olkcodebash, language=bash, caption={Running the linter}]
npm run lint
\end{lstlisting}

The config (\pathname{eslint.config.mjs}) is Next.js's flat-config defaults (\code{core-web-vitals} + \code{typescript}), plus one project-specific override: \code{react-hooks/exhaustive-deps} is configured with \code{additionalHooks: "useAsyncEffect"} so the custom hook described in the note below gets the same missing-dependency checking as a native \code{useEffect} call, rather than going unchecked because its name doesn't match what the rule looks for by default. \pathname{.github/workflows/ci.yml} runs \code{npm run lint} and \code{npm test} on every push and pull request against \code{main}, and both pass cleanly.

\begin{olknote}
\code{eslint-config-next}'s React-Compiler-era rules (\code{react-hooks/immutability}, \code{react-hooks/purity}, \code{react-hooks/set-state-in-effect}, among others) are stricter than what most of this codebase's pages were originally written against — a large number of pages used the ordinary \code{useEffect(() => \{ loadX() \}, [deps])} pattern with \code{loadX} declared below the effect, which several of these newer rules flag. Wiring up CI surfaced this pre-existing lint debt (roughly 70 errors across some twenty files); it has since been fixed, mechanically, file by file:

\begin{itemize}
  \item \textbf{"accessed before declared" (\code{react-hooks/immutability})} — the data-fetching function is now declared \emph{before} the \code{useEffect} that calls it, everywhere. The function's own hoisting doesn't matter to this rule; only textual order does.
  \item \textbf{"setState synchronously within an effect" (\code{react-hooks/set-state-in-effect})} — once a function is reachable from the effect (per the reordering above), the linter traces into it and flags any setState call reachable without an intervening callback boundary — including the extremely common \code{setLoading(true)} as an async function's first line. Wrapping the effect's call site in \code{queueMicrotask(() => fn())} (or, where one already existed, a \code{setTimeout}) breaks that trace without changing observable behavior — the callback still runs effectively immediately, before the next paint. Genuine "derive state from a changed prop" cases (\pathname{dashboard-shell.tsx}'s sidebar-close-on-navigation) were instead rewritten to adjust state directly during render, comparing against a tracked previous value — the pattern react.dev itself recommends over an effect for that specific shape.
  \item \textbf{"impure function during render" (\code{react-hooks/purity})} — \code{Date.now()}/\code{Math.random()} calls are flagged wherever they appear \emph{lexically inside} a component or hook body, regardless of whether that code path actually runs during render (e.g. inside an upload handler) — but \emph{not} inside a plain function declared outside the component and merely called from within it. The fix already used elsewhere in this codebase (\pathname{requests/page.tsx}'s top-level \code{dueDaysLeft} helper, now shared via \pathname{src/lib/date-utils.ts}) generalizes: extract the impure call into a small top-level helper.
  \item \textbf{A false-positive \code{react-hooks/rules-of-hooks}} fired on \pathname{admin/notifications/page.tsx}'s \code{useTemplate} function, because its name merely \emph{started with} "use" — it was an ordinary function, not a hook. Renamed to \code{applyTemplate}.
\end{itemize}

The remaining lint output — 37 warnings, \code{react-hooks/exhaustive-deps} and \code{@next/next/no-img-element} — has since been closed out too; \code{npm run lint} is now silent (Chapter~\ref{ch:known-issues}).
\end{olknote}

\begin{olknote}
The \code{queueMicrotask(() => fn())} fix above was applied file by file, by hand, in roughly twenty places — leaving \code{useEffect(() => \{ queueMicrotask(() => fn()) \}, deps)} duplicated verbatim across every one of them. A later code-quality pass consolidated the identical seventeen of those call sites into a shared hook, \code{useAsyncEffect(fn, deps)} (\pathname{src/hooks/use-async-effect.ts}), that wraps exactly that pattern once. The remaining handful of call sites were deliberately \emph{not} folded in: \pathname{browse/page.tsx} and \pathname{clubs/page.tsx} each pair a mount-only effect (now using \code{useAsyncEffect}) with a second, differently-shaped effect that skips its first invocation via a \code{mounted} ref before refetching on a specific dependency change — a distinct pattern, not the same one, so it was left as a direct \code{useEffect}/\code{queueMicrotask} call rather than forced into an abstraction that didn't quite fit. \pathname{messages/[id]/page.tsx}'s effect mixes the fetch with a realtime subscription setup in the same effect body, for the same reason.
\end{olknote}

\section{Automated tests: a minimal but real foundation}

\pathname{vitest.config.ts} configures Vitest with the same \code{@/*} alias as \pathname{tsconfig.json}; \code{npm test} runs every \pathname{src/**/*.test.ts} file. Coverage today is intentionally narrow — pure, isolable logic only (\pathname{html-escape.test.ts}, \pathname{text-limits.test.ts}, \pathname{geo.test.ts}) — not component tests, integration tests, or anything that touches Supabase. There is no Playwright or similar end-to-end runner. Until that changes, the practical discipline for anything not covered by a unit test is unchanged:

\begin{itemize}
  \item Run \code{npm run dev} against the \emph{local} Supabase stack, never the hosted one, while iterating.
  \item After any schema change, run \code{supabase db reset} and re-exercise the affected page end to end before considering the change done.
  \item After any RLS or admin-role change, test as more than one role — a regular user \emph{and} at least a \code{moderator} account — since RLS bugs are invisible if you only ever test as a \code{super\_admin}.
  \item When you add a new pure function (a formatter, a validator, a small parser), add a test file next to it rather than only exercising it by hand — that's how the three files above got there, and it's the cheapest way to grow real coverage incrementally.
\end{itemize}

\section{Git conventions}

Recent commit history is the best guide to house style — short, imperative, present-tense summaries of user-visible change (\emph{"Redesign the requests page as clear per-exchange cards,"} \emph{"Add distance radius filter to Browse"}). Match that tone: describe the outcome, not the mechanism. There is no enforced branch-naming or PR-template convention as of this edition.

\begin{olktask}
Create a migration with \code{npx supabase migration new my\_test\_migration} that adds a throwaway column to a table you don't care about (e.g. a \code{test\_note text} column on \code{genres}), run \code{supabase db reset}, confirm the column exists via Studio, then write and run a second migration that drops it again. This is the full round-trip you'll use for every real schema change from here on.
\end{olktask}
